\documentclass[11pt]{article}

\usepackage[margin=1in]{geometry}
\usepackage{amsmath,amssymb,amsthm}
\usepackage{mathtools}
\usepackage{enumitem}
\usepackage{booktabs}
\usepackage[colorlinks=true,linkcolor=blue,citecolor=blue,urlcolor=blue]{hyperref}

\newcommand{\GF}{\mathbb{F}_2}
\newcommand{\HX}{H_X}
\newcommand{\HZ}{H_Z}
\newcommand{\rank}{\operatorname{rank}}
\newcommand{\wt}{\operatorname{wt}}
\newcommand{\supp}{\operatorname{supp}}
\newcommand{\diam}{\operatorname{diam}}
\newcommand{\eff}{kd^2/n}

\title{\textbf{The QEC Challenge}\\[4pt]
	\large A public, automatically verified leaderboard for quantum
	low-density parity-check codes}
\author{Unitary Foundation \and
	\texttt{github.com/unitaryfoundation/qldpc-challenge}}
\date{\vspace{-2ex}\today}

\begin{document}
	\maketitle

	\begin{abstract}
		The QEC Challenge is a public leaderboard for quantum low-density
		parity-check codes. A participant submits a single JSON file describing
		one CSS code; a verifier recomputes the code's properties and, if it
		holds up, the code is merged and appears on the board. Three principles
		shape the design. (i) \emph{Everything cheap is trustless and
		mandatory}: the verifier recomputes $n$, $k$, CSS commutation, check
		weight, and geometric locality from scratch over $\GF$ --- a wrong
		claimed $k$ or a non-commuting ``stabilizer'' is rejected in
		milliseconds. (ii) \emph{Distance is split into a self-certifying upper
		bound and a separately earned exact tier}: because computing distance
		is NP-hard, a submitter attaches an explicit logical-operator
		\emph{witness} that proves $d\le\text{value}$ with no trust required;
		upgrading a claim to \emph{exact} requires a separate, more expensive
		server certification. (iii) \emph{A code is not a single number}: there
		is no one winner but a set of \emph{tracks} (hard-constraint
		categories), and within each track the ranking is a Pareto frontier
		over $(n,k,d)$ rather than a scalar. The whole pipeline ---
		validator, distance certifier, scorer, static-site generator --- lives
		in the repository and runs in CI on every pull request.
	\end{abstract}

	\section{Overview}

	A submission is one JSON file placed under \texttt{codes/}, conforming to
	\texttt{schema/code.schema.json} (described in \texttt{schema/SCHEMA.md}). A
	contributor opens a pull request adding the file; GitHub Actions runs the
	verifier (\texttt{verify/qldpc\_verify.py}), and a green check means the
	code's cheap, trustless properties are confirmed. After merge,
	\texttt{site/build.py} regenerates the static leaderboard into
	\texttt{docs/} (an index with per-track Pareto plots and tables, a page per
	code, a reference list, and status badges). The boards are seeded with the
	planar codes of Liang, Eberhardt, and Chen (arXiv:2504.08887) and other
	literature baselines, so a new entry is always measured against the
	published state of the art rather than an empty board.

	Two things are deliberately \emph{not} asked of the submitter: a proof that
	their distance is exact (that is the server's job), and a single headline
	rank (the board is a Pareto frontier). Everything else --- the parity
	checks, the claimed $n,k,d$, an optional layout, and provenance --- is
	supplied and then recomputed or witnessed by the verifier.

	\section{Why a leaderboard of tracks, not one number}

	A quantum code trades several quantities against one another --- physical
	qubits $n$, logical qubits $k$, distance $d$, maximum check weight, and
	geometric locality --- and separately has a decoding performance. Collapsing
	all of that into one scalar would reward gaming one axis at the expense of
	the rest. So the board is stratified:

	\paragraph{Tracks} are hard-constraint categories. A code only appears on a
	track's board if it satisfies that track's constraints, which the verifier
	checks. The tracks live in \texttt{TRACKS.md}; the initial set is:
	\begin{itemize}[nosep]
		\item \texttt{weight-6} --- CSS codes with all checks of weight
		$\le 6$, any connectivity (the headline frontier), with
		\texttt{weight-4} and \texttt{weight-8} as sub-/super-threshold
		siblings;
		\item \texttt{2d-local-bilayer} --- geometrically 2D-local on up to two
		physical layers (the flip-chip regime), requiring a \texttt{locality}
		block with \texttt{layers}~$\le 2$, ranked within a maximum interaction
		radius;
		\item \texttt{2d-local-single} --- 2D-local on a single layer
		(surface-code-like connectivity); stricter, mostly a baseline track;
		\item \texttt{bivariate bicycle (periodic)} --- bivariate bicycle codes
		on a torus, no 2D-local layout, seeded with the canonical codes of
		Bravyi et al.\ (arXiv:2308.07915);
		\item \texttt{generalized bicycle} --- univariate quasi-cyclic codes
		from two circulant polynomials, seeded from the literature
		(arXiv:2203.17216);
		\item \texttt{topological} --- foundational surface/toric/color codes,
		exact by construction;
		\item \texttt{decoding} (planned) --- a server-run logical-error-rate /
		threshold track under a fixed circuit-level noise model, which needs
		decoder sandboxing and lands after the parameter tracks.
	\end{itemize}
	A code may enter several tracks at once (it lists them all in
	\texttt{tracks}). New tracks can be proposed by PR.

	\paragraph{Pareto frontier, not a single winner.} Within a track, a code is
	on the frontier if no other code in that track dominates it --- i.e.\ no
	other has $n'\le n$, $k'\ge k$, $d'\ge d$ with at least one strict
	inequality. There can be many co-leaders. The board also offers a
	\emph{cell view}: a grid keyed by $(k,d)$, each cell holding the smallest
	known $n$, which is the code-tables-style view people usually screenshot.

	\paragraph{$\eff$ as a sortable column, not the rank.} The
	encoding-efficiency figure $\eff$ that the 2D-locality literature uses is
	reported and sortable, and serves as a reasonable per-track headline, but
	it never collapses the frontier into one number.

	\section{Background: the $\eff$ metric and the BPT bound}

	For an $[[n,k,d]]$ stabilizer code, the Bravyi--Poulin--Terhal (BPT) bound
	states that any geometrically 2D-local code obeys
	\begin{equation}
		\frac{kd^2}{n} \;\le\; c(r,w,q),
		\label{eq:bpt}
	\end{equation}
	where the constant $c$ depends only on the interaction range $r$, the
	maximum stabilizer weight $w$, and the on-site dimension $q$ --- not on $n$.
	The optimal constant is open, which makes ``within a fixed locality model,
	maximize $\eff$'' a well-posed target for the 2D-local tracks. The bound is
	asymptotic, so small finite codes routinely exceed any limiting family's
	constant; this is exactly why the board buckets and compares like with like
	rather than ranking a $[[120,\dots]]$ entry against a $[[10^4,\dots]]$ one.
	Reference points from the literature --- the published planar records
	(Liang--Eberhardt--Chen) alongside the periodic and topological baselines
	the board is seeded against (all $q=2$; $w$ = max check weight):
	\begin{center}
		\begin{tabular}{lcccc}
			\toprule
			Code & $[[n,k,d]]$ & topology & $w$ & $\eff$\\
			\midrule
			Toric (Kitaev), $L\times L$  & $[[2L^2,2,L]]$  & periodic & 4 & $1$\\
			Rotated planar surface       & $[[d^2,1,d]]$   & \textbf{planar} & 4 & $1$\\
			Gross code (IBM)             & $[[144,12,12]]$ & periodic & 6 & $12$\\
			Liang--Eberhardt--Chen       & $[[450,11,15]]$ & \textbf{planar} & 6 & $5.50$\\
			Liang--Eberhardt--Chen       & $[[282,12,14]]$ & \textbf{planar} & 8 & $8.34$\\
			\bottomrule
		\end{tabular}
	\end{center}
	The planar weight-6/8 state of the art is set by Liang, Eberhardt, and Chen
	(arXiv:2504.08887), who convert bivariate-bicycle codes to open boundaries
	via boundary anyon condensation plus a ``lattice grafting'' optimization,
	with all reported distances computed exactly by integer programming ---
	exactly the referee posture this challenge adopts. Whether the gap to the
	periodic weight-6 records is intrinsic to open boundaries or an artifact of
	a restricted search is the kind of question a live leaderboard can resolve.

	\section{Submission format}

	A submission is one JSON object (\texttt{schema\_version} \texttt{"0.1"},
	\texttt{code\_type} \texttt{"CSS"} --- the only type in v0.1) with:
	\begin{itemize}[nosep]
		\item \texttt{n} --- physical qubit count; must match the indices used
		in \texttt{checks}.
		\item \texttt{k} --- \emph{claimed} logical qubit count; the verifier
		recomputes it and requires an exact match.
		\item \texttt{checks.X}, \texttt{checks.Z} --- the parity checks as
		sparse supports: each is a list of checks, each check a sorted list of
		distinct $0$-based qubit indices. Thus $\HX$ has \texttt{len(checks.X)}
		rows.
		\item \texttt{distance.d} --- claimed distance $=\min(d_X,d_Z)$, and per
		side \texttt{distance.X}/\texttt{distance.Z}, each a
		$\{\texttt{value},\texttt{confidence},\texttt{witness}\}$ triple, where
		\texttt{confidence} is \texttt{"upper\_bound"} or \texttt{"exact"} and
		\texttt{witness} is the support of a logical operator of that Pauli type
		and weight \texttt{value}.
		\item \texttt{locality} (optional; required for the 2D-local tracks) ---
		\texttt{coordinates} (one $[x,y]$ per qubit), \texttt{layers}, and a
		claimed \texttt{interaction\_radius} (recomputed by the verifier).
		\item \texttt{provenance} --- \texttt{authors}, \texttt{construction},
		optional \texttt{references} (arXiv id / DOI resolved against
		\texttt{refs.bib}), \texttt{date}, \texttt{notes}.
		\item \texttt{tracks} --- the track ids this code is entered into.
	\end{itemize}
	Verify locally before opening the PR (the project uses \texttt{uv}):
	\begin{center}
		\texttt{uv run python verify/qldpc\_verify.py codes/your-code.json}
	\end{center}
	The verifier prints a JSON report --- per-check pass/fail, the computed
	$n,k,\text{ranks},\text{max\_check\_weight},\text{interaction\_radius}$, and
	an \texttt{earned\_distance} block giving the tier each side actually earned
	--- and exits $0$ iff every required check passes.

	\section{What the verifier checks}
	\label{sec:checks}

	A submission is rejected unless all of the following hold. Checks
	(V1)--(V7) are polynomial-time and run in CI on every PR; the exact
	distance tier (Section~\ref{sec:distance}) is a separate server step.
	\begin{enumerate}[label=(V\arabic*),leftmargin=*,nosep]
		\item \textbf{Schema and indices.} The document conforms to
		\texttt{code.schema.json}, and every qubit index in \texttt{checks} is
		in $[0,n)$.
		\item \textbf{No collapsed checks.} No qubit index is repeated within a
		single check (a repeat would XOR away and silently change the code).
		\item \textbf{CSS commutation.} $\HX\HZ^{\!\top}=0$ over $\GF$.
		\item \textbf{Logical dimension.} The verifier recomputes
		$k = n-\rank_{\GF}(\HX)-\rank_{\GF}(\HZ)$ and requires it to equal the
		claimed $k$.
		\item \textbf{Distance witnesses.} For each provided side, the witness
		$v$ must (a) have weight equal to the claimed \texttt{value}, (b) lie in
		the kernel of the opposite checks (an $X$-witness commutes with $\HZ$),
		and (c) be nontrivial, i.e.\ not in the rowspace of its own checks (not
		a stabilizer product). This certifies $d_{\text{side}}\le\texttt{value}$
		as an \emph{upper bound} with no trust required.
		\item \textbf{Distance consistency.} The claimed $d$ equals the minimum
		over the witnessed sides.
		\item \textbf{Locality (if present).} \texttt{coordinates} cover all $n$
		qubits, and the recomputed maximum check diameter is $\le$ the claimed
		\texttt{interaction\_radius}.
	\end{enumerate}
	Every \texttt{exact} confidence claim is \emph{downgraded to
	upper\_bound here} and flagged for server certification --- the cheap PR
	check never upgrades a distance to exact.

	\section{Computing and certifying the distance}
	\label{sec:distance}

	The distance is the minimum weight over nontrivial logicals; for a CSS code
	$d=\min(d_X,d_Z)$. The problem is NP-hard in general, and the hardness is
	asymmetric: exhibiting a low-weight logical gives only an \emph{upper}
	bound, while certifying $d\ge D$ is the hard direction, and a submitter has
	every incentive to overclaim. The challenge handles this with two tiers.

	\paragraph{Upper-bound tier (trustless, in CI).} The submitter's witness is
	checked exactly as in (V5). A valid witness proves $d_{\text{side}}\le
	\texttt{value}$; the board labels such a code $d\le$ and lets anyone climb
	instantly, since the arithmetic is checked rather than trusted.

	\paragraph{Exact tier (server-certified).} To upgrade a side to $d=$, a
	maintainer runs \texttt{verify/certify.py}, which proves no lighter
	nontrivial logical exists via a bounded integer program (scipy/HiGHS, no
	commercial solver). For each logical generator $t_L$ of the relevant type it
	asks whether a $v$ exists with
	\begin{equation}
		H v \equiv 0,\qquad \langle v,t_L\rangle\equiv 1 \pmod 2,\qquad
		\wt(v)\le \texttt{value}-1,
		\label{eq:milp}
	\end{equation}
	the mod-2 constraints linearized by integer slacks. If every generator on
	both sides is proven \emph{infeasible}, $d$ is exact; if a lighter logical
	is found, the claim is \emph{refuted}; if the solver hits its time limit, the
	side honestly remains an upper bound. A successful run writes a certificate
	to \texttt{certs/<slug>.json} (with \texttt{d\_exact}, per-side notes, solver,
	and per-solve time limit), which is what upgrades a code's board label to
	$d=$. Exact records outrank equal upper-bound ones; nothing is ever silently
	upgraded.

	\section{Duplicate detection}

	To keep the board honest about novelty, the verifier fingerprints every
	code two ways. The reduced-row-echelon forms of $\HX,\HZ$ pin the exact
	stabilizer group: two codes with the same RREF fingerprint are
	\emph{identical} (same labeling) and a duplicate is a hard CI error. A
	finer, permutation-invariant Weisfeiler--Leman color refinement on the
	Tanner graph gives a signature that any two qubit-permuted copies must
	share; a collision \emph{flags} a possible equivalent for human review (WL
	is a strong necessary condition, not a complete equivalence test).

	\section{Repository layout and automation}

	\begin{center}
		\begin{tabular}{ll}
			\toprule
			Path & Contents\\
			\midrule
			\texttt{schema/}   & submission format (JSON Schema + \texttt{SCHEMA.md})\\
			\texttt{verify/}   & verifier: \texttt{gf2.py} (GF(2) core),
			                     \texttt{qldpc\_verify.py} (checker),\\
			                  & \texttt{certify.py}/\texttt{certify\_all.py} (exact tier),
			                     \texttt{verify\_all.py}, tests\\
			\texttt{examples/} & worked examples that pass verification\\
			\texttt{codes/}    & accepted submissions (the leaderboard data)\\
			\texttt{certs/}    & server-side exact-distance certificates\\
			\texttt{site/}     & \texttt{build.py}, which generates the static site\\
			\texttt{docs/}     & the generated site (index, per-code pages,
			                     references, badges)\\
			\texttt{refs.bib}  & bibliography; \texttt{references.html} is generated from it\\
			\texttt{TRACKS.md} & the tracks and how ranking works\\
			\bottomrule
		\end{tabular}
	\end{center}
	CI (\texttt{.github/workflows/verify.yml}) runs the verifier self-tests
	(\texttt{test\_verifier.py}, which must reject malformed and dishonest
	submissions) and then \texttt{verify\_all.py} over every file in
	\texttt{codes/} and \texttt{examples/}, failing if any code is invalid or if
	two entries share an identical fingerprint. The badges in the README and the
	figures in \texttt{docs/stats.json} are regenerated on every build from the
	live board, so they track the current numbers rather than a hand-typed
	count. At the time of writing the board holds on the order of three dozen
	verified codes (around twenty certified exact) across six populated tracks,
	with a best $\eff\approx 14.2$ (the $[[126,28,8]]$ generalized-bicycle code
	of Panteleev--Kalachev, weight-10 checks); the live figures are in
	\texttt{docs/stats.json}.

	\section{Open questions this could surface}

	\begin{itemize}[nosep]
		\item Can planar weight-6 codes close the gap to the periodic records,
		or is part of it intrinsic to open boundaries?
		\item Do bulk stabilizers or boundary geometries outside the families
		searched so far beat the Liang et al.\ records?
		\item The small-lattice distances of the seed codes are set by
		Sierpinski-type \emph{fractal} logical operators rather than strings ---
		can this regime be engineered deliberately for better finite-size codes?
	\end{itemize}

	\subsection*{Key references}
	\begin{itemize}[nosep]
		\item Liang, Eberhardt, Chen, arXiv:2504.08887 --- planar open-boundary
		qLDPC codes: the seed planar baseline (anyon-condensation $+$
		lattice-grafting construction).
		\item Bravyi, Cross, Gambetta, Maslov, Rall, Yoder, Nature
		\textbf{627}, 778 (2024), arXiv:2308.07915 --- the $[[144,12,12]]$ gross
		code and bivariate-bicycle codes (periodic-track seed).
		\item Wang, Pryadko, Symmetry \textbf{14}, 1348 (2022),
		arXiv:2203.17216 --- distance bounds for generalized bicycle codes; and
		Panteleev, Kalachev, Quantum \textbf{5}, 585 (2021), arXiv:1904.02703
		--- generalized bicycle codes including the $[[126,28,8]]$ headline
		entry (generalized-bicycle-track seeds).
		\item Kitaev, Ann.\ Phys.\ \textbf{303}, 2 (2003),
		arXiv:quant-ph/9707021 --- the toric code; Steane,
		arXiv:quant-ph/9601029 --- CSS codes (topological-track baseline).
		\item Panteleev, Kalachev, STOC 2022, arXiv:2111.03654, and Tillich,
		Z\'emor, IEEE Trans.\ IT \textbf{60}, 1193 (2014), arXiv:0903.0566 ---
		good and hypergraph-product qLDPC codes.
		\item Pryadko, Shabashov, Kozin, \texttt{QDistRnd}, JOSS \textbf{7}, 4120
		(2022) --- randomized distance estimation, usable as a refutation
		oracle; Perlin, \texttt{qLDPC} (\url{https://github.com/qLDPCOrg/qLDPC})
		--- a convenient way to construct codes and export their parity checks.
	\end{itemize}

\end{document}
